\chapter{Analyse des besoins et spécifications fonctionnelles}

\section{Démarche de recueil et d'analyse des exigences}

L'élicitation des exigences du système École~221 a été menée selon une démarche rigoureuse d'ingénierie logicielle basée sur l'analyse documentaire approfondie des règlements scolaires et la modélisation des processus d'affaires des établissements d'enseignement supérieur et secondaire.

Conformément aux directives de transparence, les exigences ont été consolidées dans le document formel (\nolinkurl{docs/requirements.md}). Chaque exigence porte l'identifiant \texttt{US-xx} pour les récits utilisateurs et \texttt{BR-xx} pour les règles métier, avec des critères d'acceptation testables par automate.

\section{Acteurs du système et cas d'utilisation}

Le système s'articule autour de six profils d'utilisateurs distincts régis par un contrôle d'accès basé sur les rôles (\emph{Role-Based Access Control - RBAC}) :

\begin{itemize}
  \item \textbf{Super-Admin (SA)} : gestion globale des établissements (tenants), allocation des quotas d'utilisateurs, suivi de l'audit système et paramétrage des passerelles techniques ;
  \item \textbf{Directeur / Secrétariat (DIR)} : administration générale de l'école, ouverture des campagnes d'admission, validation des dossiers et affectation des enseignants ;
  \item \textbf{Comptable / Gestionnaire financier (COMPTA)} : émission des factures de scolarité, encaissement des versements, suivi du tableau de bord financier ;
  \item \textbf{Enseignant (PROF)} : saisie des présences, création des devoirs et saisie des notes d'évaluations ;
  \item \textbf{Parent / Tuteur légal (PAR)} : consultation du bulletin de notes en ligne, justification des absences et paiement des frais d'inscription ;
  \item \textbf{Vigile / Agent de sécurité (VIG)} : contrôle d'accès aux portes via scanner QR Code avec fonctionnement garanti en mode hors-ligne.
\end{itemize}

\section{Exigences fonctionnelles prioritaires}

Le tableau synthétique ci-dessous récapitule les principales exigences fonctionnelles implémentées dans le prototype validé :

\begin{table}[htbp]
  \centering
  \small
  \begin{tabular}{lllp{6.5cm}c}
    \toprule
    \textbf{ID} & \textbf{Domaine} & \textbf{Acteur} & \textbf{Description du besoin} & \textbf{Priorité} \\
    \midrule
    \texttt{US-01} & Multi-Tenancy & Super-Admin & Créer un établissement scolaire avec son sous-domaine et ses administrateurs dédiés. & P0 (Vitale) \\
    \texttt{US-02} & Inscription & Candidat & Soumettre un dossier de candidature complet avec pièces jointes justificatives. & P0 (Vitale) \\
    \texttt{US-03} & Pédagogie & Secrétariat & Configurer les filières, classes, semestres et coefficients des matières. & P1 (Haute) \\
    \texttt{US-04} & Évaluation & Enseignant & Saisir les notes de contrôle continu et d'examen avec calcul automatisé des moyennes. & P0 (Vitale) \\
    \texttt{US-05} & Finance & Comptable & Émettre l'échéancier de scolarité et enregistrer les règlements d'élèves. & P1 (Haute) \\
    \texttt{US-06} & Sécurité & Vigile & Scanner un badge QR physique ou mobile pour autoriser ou refuser l'entrée d'un élève. & P0 (Vitale) \\
    \texttt{US-07} & Communication & Parent & Recevoir des alertes immédiates en cas d'absence injustifiée ou de publication de notes. & P2 (Moyenne) \\
    \bottomrule
  \end{tabular}
  \caption{Matrice des exigences fonctionnelles prioritaires de la plateforme École 221}
  \label{tab:exigences-fonctionnelles}
\end{table}

\section{Exigences non fonctionnelles et sécurité}

Les exigences non fonctionnelles ont été calibrées selon la norme internationale ISO/IEC~25010~\cite{iso25010} afin de garantir la pérennité industrielle de la solution :

\begin{itemize}
  \item \textbf{Étanchéité et cloisonnement des données (SEC-003)} : interdiction absolue de toute fuite de données inter-écoles. Le système doit garantir que chaque requête SQL filtre automatiquement par la clé d'organisation de l'école authentifiée, sans dépendre d'un paramètre transmis par le navigateur ;
  \item \textbf{Temps de réponse et performance} : 95\% des requêtes d'API REST de consultation doivent répondre en moins de 250~millisecondes sur une infrastructure Cloud standard ;
  \item \textbf{Résilience et fonctionnement dégradé} : le terminal vigile doit continuer à valider les accès au campus même lors d'une coupure complète de la liaison Internet, grâce à un cache chiffré local et une resynchronisation différée ;
  \item \textbf{Auditabilité et traçabilité} : chaque action administrative critique (modification d'une note après clôture, suppression d'un compte, bascule de statut de paiement) doit faire l'objet d'un journal d'audit immuable horodaté.
\end{itemize}
